\section{Method}
\label{sec:method}

We present a retrieval framework that operates over a structural representation of a software repository---a \emph{graph world model}---and a patching pipeline that uses retrieved context to generate unified diffs. The framework supports eleven retrieval strategies spanning three tiers: zero-LLM-runtime baselines, single-turn agentic ablations, and multi-turn agentic methods. This section describes the graph world model (\S\ref{sec:graph_world_model}), the retrieval methods in each tier (\S\ref{sec:retrieval_methods}), the patching pipeline (\S\ref{sec:patching_method}), and the cost model used for comparison (\S\ref{sec:cost_model}).

\subsection{Graph World Model}
\label{sec:graph_world_model}

Given a repository snapshot at a specific commit, we construct a directed, typed knowledge graph $G = (V, E)$ using static analysis of the Python source tree via the \texttt{tree-sitter} parser. Nodes $V$ are categorised into three types: \emph{file} nodes (one per source file), \emph{class} nodes, and \emph{function/method} nodes. Each node stores a compact metadata record comprising its qualified name, file path, docstring (if present), and a brief signature. Node embeddings are computed once at build time over this metadata using the Gemini text-embedding-004 model; the resulting vectors populate a FAISS IVF index for semantic search.

Edges $E$ encode four structural relation types: \textsc{contains} (file $\to$ class, file $\to$ function, class $\to$ method), \textsc{imports} (file $\to$ file, resolved statically), \textsc{calls} (function $\to$ function, with confidence tier: \emph{definite}, \emph{probable}, or \emph{possible} based on static resolvability), and \textsc{inherits} (class $\to$ parent class). The graph index is built once per repository snapshot; across a retrieval evaluation run on a single snapshot, the setup cost is amortised over all issues evaluated against that snapshot.

\subsection{Retrieval Methods}
\label{sec:retrieval_methods}

We evaluate eleven retrieval strategies, grouped into three tiers by their use of LLM calls at query time.

\subsubsection{Zero-LLM-Runtime Tier}
\label{sec:tier_zero}

These methods issue no LLM calls during query execution. Their total token cost consists only of the one-time graph setup embedding.

\paragraph{BM25 (\texttt{bm25}).}
BM25Plus \citep{lv2011lower} over tokenised source files provides a non-neural baseline with no index. Files are scored by term-overlap between the issue text and the concatenation of file path, function names, docstrings, and source tokens. No graph structure or embeddings are used.

\paragraph{Graph-PageRank (\texttt{repomap\_like}).}
Loosely inspired by Aider's repository map concept \citep{aider2023}, this method computes a personalised PageRank over the file-level import graph, treating files that have high BM25 similarity to the issue as seed nodes. The final file score is the product of the BM25 score and the PageRank mass, biasing retrieval toward structurally central files that are lexically relevant to the issue. \textbf{Note:} unlike Aider's original mechanism, which uses a compressed AST-based text representation of the repository, our adaptation replaces this with personalised PageRank over the static import graph. This is a controlled architectural comparison, not a faithful reproduction of the Aider system.

\paragraph{GM-Deterministic (\texttt{gm\_deterministic}).}
A multi-component scoring function over all nodes in the graph world model, requiring no LLM calls at query time. For each node, a score is computed as a weighted sum of four components: (1) \emph{semantic score}---cosine similarity between the query embedding and the node embedding; (2) \emph{structural score}---a BM25 match of the issue text against the node's graph neighbourhood (immediate import/call neighbours); (3) \emph{confidence score}---inversely proportional to the minimum call-edge confidence tier on paths from the query seed to the node; and (4) \emph{hub penalty}---a multiplicative discount applied to nodes with very high in-degree (hub nodes that appear spuriously relevant due to centrality). The top-$K$ scored files are returned.

\paragraph{Raw-RAG-Function (\texttt{raw\_rag\_function}).}
Dense retrieval over function-body embeddings. Each function (or method) in the repository is embedded separately using its full source text. At query time, the top-$K$ most similar function embeddings are retrieved and their containing files are returned. This establishes a fine-grained dense-retrieval baseline that does not use graph structure.

\paragraph{Raw-RAG-Fixed (\texttt{raw\_rag\_fixed}).}
Dense retrieval over fixed-size sliding-window chunks. The source code of each file is chunked into overlapping windows of a fixed token budget, each chunk embedded independently. Top-$K$ chunks by cosine similarity are retrieved and their source files returned. This corresponds to the standard code-RAG approach used in production retrieval pipelines.

\subsubsection{Single-Turn Agentic Tier (Ablation Variants)}
\label{sec:tier_single}

These methods make a single LLM call at query time to either score or query the graph. They appear in the retrieval evaluation to isolate the contribution of multi-turn reasoning, but are excluded from patching runs (which require high-quality file lists).

\paragraph{GM-Baseline (\texttt{gm\_baseline}).}
A single LLM turn is permitted; the model may invoke the same graph tools as GM-Progressive (\texttt{search\_nodes}, \texttt{get\_neighbors}, \texttt{get\_file\_summary}) but receives no follow-up turns. This ablation isolates the benefit of iterative multi-turn graph navigation over single-pass tool-assisted scoring.

\paragraph{RAG-Baseline (\texttt{rag\_baseline}).}
A single LLM turn is permitted; the model may invoke one round of the same symmetric tools as RAG-Progressive (\texttt{semantic\_search}, \texttt{get\_file\_contents}, \texttt{confirm\_file}) but receives no follow-up turns. This ablation isolates the benefit of multi-turn query refinement in RAG-Progressive.

\subsubsection{Multi-Turn Agentic Tier}
\label{sec:tier_multi}

These methods engage the LLM in an iterative dialogue with retrieval tools, allowing progressive refinement.

\paragraph{GM-Progressive (\texttt{gm\_progressive}).}
The primary graph-guided method. The LLM is provided with a tool set for graph traversal: \texttt{search\_nodes} (semantic search over the FAISS index), \texttt{get\_neighbors} (expand a node's immediate graph neighbourhood), and \texttt{get\_file\_summary} (list all definitions in a file and implicitly mark it as confirmed). In each turn, the model may call any tool or emit a final answer. The dialogue continues until the model emits a final file list or a stagnation limit is reached (consecutive turns with no new file confirmations). The multi-turn structure allows the model to progressively build confidence in which files are relevant, using graph structure to resolve structural dependencies that embedding similarity alone cannot capture.

\paragraph{RAG-Progressive (\texttt{rag\_progressive}).}
Multi-turn RAG using symmetric tools: \texttt{semantic\_search} (dense vector search), \texttt{get\_file\_contents} (read a source file), and \texttt{confirm\_file}. The model may issue multiple queries with refined terms and inspect file contents before committing to a final list. This provides a strong agentic baseline that benefits from multi-turn refinement without graph structure.

\paragraph{Agentless-Like (\texttt{agentless\_like\_localization}).}
A three-stage hierarchical localization procedure adapted from Agentless \citep{xia2024agentless}. Stage~1: the LLM receives the directory tree and ranks candidate files by relevance to the issue. Stage~2: the LLM receives graph-derived symbol-level structure (functions, classes, and their signatures) for the Stage~1 candidates and selects relevant symbols. Stage~3: the LLM receives the selected symbol bodies and identifies the precise relevant spans, which are then mapped back to file-level granularity to conform to the unified file-list interface used across all methods. \textbf{Important:} unlike the original Agentless, our implementation uses the GraphManager dependency graph to supply symbol-level structure in Stage~2, making this a partially graph-augmented method. Additionally, our implementation generates a single repair candidate, whereas the original Agentless uses 40 samples with execution-based filtering. Results for this method should be interpreted with these differences in mind (see Section~\ref{sec:threats}).

\paragraph{Agentic Cold-Start (\texttt{agentic\_cold\_start}).}
Direct filesystem tool use with no retrieval index. The LLM is provided with tools to navigate and read the raw repository: \texttt{list\_directory}, \texttt{read\_file}, and \texttt{search\_code} (regex over file contents). No graph, no embedding index. This method serves as a baseline for the intrinsic file-navigation capability of the model and as an upper bound on what is achievable without retrieval infrastructure.

\subsection{Patching Method}
\label{sec:patching_method}

Patch generation is a single-turn procedure applied uniformly across all retrieval methods that produce a confirmed file list. The issue description and the full contents of the confirmed files are concatenated into a single context window, and a Gemini 2.0 Flash model generates a unified diff in standard \texttt{git diff} format. No multi-turn repair is attempted. This design isolates the effect of retrieval quality on patch success: differences in resolved rate across retrieval methods reflect retrieval quality, not patching strategy variation.

Patch evaluation uses the SWE-bench Verified harness \citep{swebench2024}: the generated patch is applied to the repository at the issue's \texttt{base\_commit} and the repository's test suite is executed in an isolated Docker container. A patch is counted as resolved if and only if all tests that were failing before the patch now pass, with no regression in previously passing tests.

\subsection{Cost Model}
\label{sec:cost_model}

Token cost is decomposed into three components for every method:

\begin{description}[leftmargin=1em,style=nextline]
  \item[Setup embedding tokens] Tokens consumed building the retrieval index (graph metadata embeddings, or code chunk embeddings for RAG methods). Amortised over all issues evaluated against a single snapshot.
  \item[Runtime LLM tokens] Prompt + completion tokens consumed by LLM calls during query execution. Zero for the Zero-LLM-Runtime tier.
  \item[Query embedding tokens] Tokens consumed embedding the issue query for dense search. Zero for BM25 and Graph-PageRank.
\end{description}

We report two cost-per-resolved-issue (CPR) views: \emph{method-accounted CPR} adds each method's full setup cost to its runtime cost before dividing by resolved count, providing a like-for-like comparison; \emph{as-run CPR} reports only the tokens actually consumed in the patching run (i.e., excluding setup costs shared with retrieval evaluation), reflecting practical amortised cost when the same index is reused.
